\chapter{Solution Overview}
\label{chap:solution}

\section{Target Platform}
\label{sec:target-platform}

Power traces are measured on a ChipWhisperer-Lite board with an
STM32F303 (Arm Cortex-M4) target, controlled over the SimpleSerial
protocol. A hardware trigger raised by the firmware just before the
forward pass delimits the neural-network inference. Each trace contains
$24{,}400$ samples, but only its first $1{,}000$ samples, which contain
the targeted multiplication, are used for the analysis. This speeds up
profiling.

The target runs a feed-forward neural network with four layers of
$7$, $5$, $4$ and $3$ neurons (ReLU in the hidden layers, sigmoid at the
output), zero biases, and all values stored as IEEE-754 single-precision
numbers. The first multiplication of the inference is
$v=\mathrm{float32}(w^\star x)$, where $x$ is the input chosen by the
attacker and $w^\star$ is the secret first-layer weight. The firmware is
compiled with the default \texttt{-mfloat-abi=soft}, so every
multiplication runs as a software routine that processes the sign,
exponent and fraction bits with integer instructions, even though the
MCU has a hardware FPU.

\section{Leakage Model}
\label{sec:leakage-model}

A single-precision number~\cite{IEEE754} consists of a sign bit, $8$ exponent bits and
$23$ fraction bits. For a candidate weight $w$ and known input $x$ we
model the leakage of the intermediate $v=\mathrm{float32}(wx)$ by the
Hamming weight of its $m$ most significant bits $B_m(v)$:
\begin{equation}
	c_m(v)=\mathrm{HW}\left(B_m(v)\right).
\end{equation}
Three resolutions are used: coarse ($m=17$, $18$ classes), fine
($m=23$, $24$ classes) and exact ($m=32$, $33$ classes). Fewer modeled
bits tolerate small errors in the candidate weight, while more bits make
the model more selective.

\section{Profiled Attack}
\label{sec:attack}

\paragraph{Profiling.}
We capture $50{,}000$ traces with weights drawn uniformly from $[-2,2]$
(refreshed every $100$ traces) and random inputs. Since $x$ and the
targeted weight are known, $v$ and its class $c_m(v)$ are known for every
trace. Five points of interest (POIs) are selected as the samples with
the highest absolute Pearson correlation with $\mathrm{HW}(B_{32}(v))$,
at least $20$ samples apart. For each class $c$ we estimate the mean
$\boldsymbol{\mu}_c$ of the POI vectors $\mathbf{t}$ and a single
pooled covariance matrix $\boldsymbol{\Sigma}$, which gives the Gaussian
templates
\begin{equation}
	\mathbf{t}\mid c\sim\mathcal{N}\left(\boldsymbol{\mu}_c,\boldsymbol{\Sigma}\right).
\end{equation}
One template set is built for each of the three resolutions.

\paragraph{Attack.}
On the target the weight is fixed and unknown, and only the input $x_i$
is controlled ($5{,}000$ attack traces are available). For a candidate
$\hat{w}$ the hypothetical class $c_i(\hat{w})$ of
$\mathrm{float32}(\hat{w}x_i)$ selects the template, and the evidence is
accumulated over the attack traces:
\begin{equation}
	L(\hat{w})=\sum_{i=1}^{N_{\mathrm{attack}}}\log p\left(\mathbf{t}_i\mid c_i(\hat{w})\right).
	\label{eq:likelihood}
\end{equation}
The candidate with the highest $L(\hat{w})$ is the recovered weight.

\section{Hierarchical Weight Recovery}
\label{sec:hierarchical-recovery}

Evaluating all $\approx 2^{32}$ \texttt{float32} values is impractical,
so the search narrows the candidates in three stages
(Table~\ref{tab:hierarchical-recovery}), each scored with
Eq.~\eqref{eq:likelihood} and a finer leakage model.

\begin{table}[htb]
	\centering
	\small
	\caption{Stages of the hierarchical weight recovery.}
	\label{tab:hierarchical-recovery}
	\begin{tabular}{@{}lllll@{}}
		\toprule
		Stage  & Search range               & Step      & Candidates         & Bits (classes) \\
		\midrule
		Coarse & $[-2,2]$                   & $10^{-4}$ & $40{,}001$         & $17$ ($18$)    \\
		Fine   & $\pm 10^{-4}$ per finalist & $10^{-6}$ & $201$ per finalist & $23$ ($24$)    \\
		Exact  & $\pm 512$ ULPs             & 1 ULP     & $\approx 22{,}550$ & $32$ ($33$)    \\
		\bottomrule
	\end{tabular}
\end{table}

The coarse stage keeps the six best peaks, at least $0.005$ apart, as
finalists. The fine stage refines each of them and passes the best one
on, and the exact stage selects the final \texttt{float32} value.
Because $\pm\hat{w}\,2^j$ have the same fraction bits as $\hat{w}$,
their products differ from the true one only in the sign and exponent
bits, and the coarse and fine models cannot reliably separate them. The
exact stage therefore also tries both signs and the scalings $2^j$,
$j\in\{-10,\ldots,10\}$, within $[-2,2]$. The experimental results are
given in Chapter~\ref{chap:scientific}.

\section{Justification of the Chosen Approach}
\label{sec:justification}

\paragraph{Template attack.}
Previous parameter-recovery attacks on neural networks are mostly
non-profiled and based on CPA or CEMA~\cite{batina2019csi,
	joud2022practical, nozaki2024power}: they correlate one hypothetical
leakage value with each sample of the trace. Template attacks were so
far applied to neural networks mainly for input
recovery~\cite{wei2018iknow}. The targeted product can take up to $33$
Hamming-weight classes, and its software implementation spreads the
leakage over several samples of the trace. A template attack models
these samples jointly together with their noise correlation and ranks
candidates by their likelihood, which suits testing a large number of
$32$-bit hypotheses. Profiling is realistic because the attacker can
program an identical device with chosen weights
(Section~\ref{sec:attack}).

\paragraph{Hierarchical search.}
The coarse model uses only the sign, exponent and upper fraction bits,
so a candidate close to the true weight already predicts mostly the
right class. The likelihood is therefore smooth enough to find promising
regions on a coarse grid, and the later stages spend more candidates and
a finer model only where needed. This avoids applying the most detailed
model to the whole candidate space.

\paragraph{Pooled covariance and few POIs.}
Many of the $33$ classes are rare in the profiling data, so per-class
covariance matrices would be estimated from very few traces. A shared
covariance matrix and five POIs keep the number of estimated parameters
small and the estimation stable with $50{,}000$ profiling traces.
